\documentclass{beamer}
\usetheme{Madrid}
\usecolortheme{default}
\setbeamertemplate{footline}[frame number]

\usepackage[T1]{fontenc}
\usepackage[utf8]{inputenc}
\usepackage{amsmath,amssymb}
\usepackage{booktabs}
\usepackage{graphicx}
\usepackage{mathtools}
\usepackage{array}
\usepackage{tipa}

\DeclareUnicodeCharacter{0259}{\textschwa}
\DeclareUnicodeCharacter{018F}{\textschwa}

\graphicspath{
  {../outputs/project3/report_figures/}
  {outputs/project3/report_figures/}
}

\title{Spring 2026 Natural Language Processing (CSCI-6515)}
\subtitle{Project 3: Word Embeddings, Comparative Analysis, DL Classification, and Interactive UI}
\author{
\textbf{Kamal Ahmadov} \\ \small kahmadov24700@ada.edu.az; kamal.ahmadov@gwu.edu \\ \and
\textbf{Rufat Guliyev} \\ \small rguliyev24988@ada.edu.az; rufat.guliyev@gwu.edu
}
\institute{ADA University \& George Washington University}
\date{March 19, 2026}

\newcommand{\pp}{\operatorname{PP}}
\newcommand{\fmacro}{F1_{\text{macro}}}

\begin{document}

\begin{frame}
  \titlepage
\end{frame}

\begin{frame}{Outline}
\tableofcontents
\end{frame}

\section{Motivation and Scope}

\begin{frame}{Motivation}
\begin{itemize}
  \item Azerbaijani NLP needs stronger semantic representations than sparse token counts.
  \item We evaluate embeddings at two levels:
  \begin{itemize}
    \item Intrinsic: nearest neighbors + analogy arithmetic.
    \item Extrinsic: downstream sentiment classification.
  \end{itemize}
  \item Project 3 objective: complete Tasks 1--6 and provide an interactive UI for exploration.
\end{itemize}
\end{frame}

\begin{frame}{Datasets and Pipeline}
\textbf{Data sources}
\begin{itemize}
  \item Corpus for Tasks 1--3: \texttt{data/raw/corpus.csv}
  \item Sentiment for Tasks 4--5: \texttt{data/external/train.csv}, \texttt{test.csv}
\end{itemize}

\vspace{0.3em}
\textbf{Pipeline orchestration}
\begin{itemize}
  \item Runner: \texttt{bash scripts/run\_project3.sh}
  \item Task outputs: \texttt{outputs/project3/task1\_...task6\_...}
  \item UI: Streamlit dashboard + legacy Flask UI
\end{itemize}
\end{frame}

\section{Task 1: Dataset + Matrices}

\begin{frame}{Task 1 Method (with light math)}
\small
\textbf{Term-document matrix}
\[
X \in \mathbb{R}^{D \times V}, \quad X_{d,w} = \text{count}(w \text{ in document } d)
\]

\textbf{Word-word co-occurrence matrix}
\[
C_{i,j} = \sum_{t}\mathbf{1}[w_t=i]\mathbf{1}[w_{t+\Delta}=j], \ |\Delta|\leq W
\]

\vspace{0.4em}
\textbf{Configured thresholds}
\begin{itemize}
  \item Frequent words: count $\ge 100$
  \item Rare words: count $\le 2$
  \item Heatmap view: top terms/docs for readability
\end{itemize}
\end{frame}

\begin{frame}{Task 1 Results}
\small
\begin{table}
\centering
\begin{tabular}{lr}
\toprule
Metric & Value \\
\midrule
Documents & 31,842 \\
Sentences & 813,319 \\
Tokens & 11,905,937 \\
Vocabulary size & 586,674 \\
Frequent words ($\ge100$) & 12,436 \\
Rare words ($\le2$) & 398,599 \\
\bottomrule
\end{tabular}
\end{table}

\vspace{0.3em}
\textbf{Artifacts:} term-document sparse matrix, word frequency distribution, top-matrix CSVs, and heatmaps.
\end{frame}

\begin{frame}{Task 1 Visualizations}
\begin{columns}[T,totalwidth=\textwidth]
\column{0.5\textwidth}
\centering
\includegraphics[width=\linewidth]{task1_matrices_term_document_heatmap.png}
\\ \tiny Term-document heatmap

\column{0.5\textwidth}
\centering
\includegraphics[width=\linewidth]{task1_matrices_word_word_heatmap.png}
\\ \tiny Word-word heatmap
\end{columns}
\end{frame}

\section{Task 2: Word2Vec}

\begin{frame}{Task 2 Word2Vec Setup}
\small
\textbf{Training objective (skip-gram)}
\[
\max \sum_{t}\sum_{-k \le j \le k,\ j\neq0} \log p(w_{t+j}\mid w_t)
\]

\vspace{0.4em}
\textbf{Configuration}
\begin{itemize}
  \item Model: skip-gram (\texttt{cbow=0})
  \item Dimension: 200, Window: 5
  \item Negative sampling: 10, Min count: 5
  \item Iterations: 10
\end{itemize}

\textbf{Evaluation}
\begin{itemize}
  \item Similar words for 10 target words
  \item Analogy arithmetic (\(b-a+c\))
\end{itemize}
\end{frame}

\begin{frame}{Task 2 Results (Word2Vec)}
\small
\begin{itemize}
  \item Vocab size: 122,930, vector dim: 200
  \item Analogy hit@k: 0.6667 (3 evaluated / 5 total)
  \item Good local morphology for several targets (e.g., seher $\rightarrow$ seheri)
  \item Some noisy neighbors remain (data noise + OOV effects)
\end{itemize}

\vspace{0.3em}
\textbf{Example analogies}
\begin{itemize}
  \item Correct: \texttt{kishi:qadin :: oglan:qiz}
  \item Correct: \texttt{baki:azerbaycan :: ankara:turkiye}
  \item Incorrect: \texttt{ata:ana :: ogul:?} $\rightarrow$ \texttt{isterem}
\end{itemize}
\end{frame}

\section{Task 3: GloVe}

\begin{frame}{Task 3 GloVe Setup}
\small
\textbf{GloVe objective}
\[
J=\sum_{i,j} f(X_{ij})\left(w_i^\top \tilde{w}_j + b_i + \tilde{b}_j - \log X_{ij}\right)^2
\]

\vspace{0.4em}
\textbf{Configuration}
\begin{itemize}
  \item Dimension: 200, Window: 10
  \item Min count: 5, Iterations: 20
  \item \(x_{\max}=10\), \(\eta=0.05\), \(\alpha=0.75\)
\end{itemize}

\textbf{Same evaluation protocol as Task 2}
\end{frame}

\begin{frame}{Task 3 Results (GloVe)}
\small
\begin{itemize}
  \item Vocab size: 122,930, vector dim: 200
  \item Analogy hit@k: 0.7500 (4 evaluated / 5 total)
  \item Strong morphology neighbors also observed for many targets
  \item Some high-score noisy neighbors still appear
\end{itemize}

\vspace{0.3em}
\textbf{Example analogies}
\begin{itemize}
  \item Correct: \texttt{kishi:qadin :: oglan:qiz}
  \item Correct: \texttt{shimal:cenub :: sherq:qerb}
  \item Incorrect: \texttt{ata:ana :: ogul:?} $\rightarrow$ \texttt{isterem}
\end{itemize}
\end{frame}

\section{Task 4: Comparison}

\begin{frame}{Task 4: Word2Vec vs GloVe}
\small
\begin{table}
\centering
\begin{tabular}{lrrrr}
\toprule
Model & Hit@1 & Hit@k & Avg top-3 cosine & Coverage \\
\midrule
Word2Vec & 0.6667 & 0.6667 & 0.6619 & 0.0749 \\
GloVe & 0.7500 & 0.7500 & 0.6217 & 0.0749 \\
\bottomrule
\end{tabular}
\end{table}

\vspace{0.3em}
\textbf{Interpretation}
\begin{itemize}
  \item GloVe is better on analogy score.
  \item Word2Vec is better on local-neighbor cosine quality.
  \item Coverage is equal; downstream choice depends on objective.
\end{itemize}
\end{frame}

\section{Task 5: DL Classification}

\begin{frame}{Task 5 Experimental Design}
\small
\textbf{Architectures}
\begin{itemize}
  \item RNN, BiRNN, LSTM
\end{itemize}

\textbf{Feature families}
\begin{itemize}
  \item Count-SVD, TFIDF-SVD, PMI-SVD, Word2Vec, GloVe
\end{itemize}

\textbf{Total runs}
\begin{itemize}
  \item 15 neural runs (3x5) + baseline (LogReg + TFIDF)
\end{itemize}

\textbf{Primary metric}
\[
\fmacro = \frac{1}{C}\sum_{c=1}^{C}F1_c
\]
\end{frame}

\begin{frame}{Task 5 Results}
\small
\begin{table}
\centering
\begin{tabular}{lrrr}
\toprule
Experiment & Accuracy & Macro-F1 & Weighted-F1 \\
\midrule
\texttt{lstm\_\_pmi\_svd} & 0.8225 & 0.5684 & 0.8594 \\
\texttt{rnn\_\_pmi\_svd} & 0.8396 & 0.5678 & 0.8641 \\
\texttt{birnn\_\_pmi\_svd} & 0.8170 & 0.5600 & 0.8541 \\
\texttt{logreg\_baseline\_\_tfidf} & 0.8271 & 0.5588 & 0.8582 \\
\bottomrule
\end{tabular}
\end{table}

\vspace{0.2em}
\textbf{Best run:} \texttt{lstm\_\_pmi\_svd} by macro-F1.
\end{frame}

\begin{frame}{Task 5 Visualization}
\centering
\includegraphics[width=0.95\linewidth]{presentation_bundle_figures_task5_top10_macro_f1.png}
\end{frame}

\section{Task 6 + Extra UI}

\begin{frame}{Task 6 Deliverables}
\small
\begin{itemize}
  \item Consolidated report tables:
  \begin{itemize}
    \item \texttt{embedding\_comparison\_table.csv}
    \item \texttt{dl\_results\_table.csv}
  \end{itemize}
  \item Report note artifact: \texttt{report\_artifacts.md}
  \item Presentation bundle zip with figures/tables/summaries
\end{itemize}
\end{frame}

\begin{frame}{Extra Task: Interactive UI}
\small
\textbf{Primary UI:} Streamlit dashboard (\texttt{src/project3\_dashboard.py})
\begin{itemize}
  \item Overview and task tabs for Task 1--6 outputs
  \item Embedding playground:
  \begin{itemize}
    \item custom nearest-neighbor queries
    \item custom analogy solver
    \item vector blending and semantic map (2D)
  \end{itemize}
  \item Read-only artifact browsing (no retraining from UI)
\end{itemize}
\end{frame}

\section{Limitations and Conclusion}

\begin{frame}{Limitations and Improvement Areas}
\small
\begin{itemize}
  \item Sentiment dataset is highly imbalanced (positive dominates).
  \item Intrinsic analogy set is relatively small and OOV-sensitive.
  \item Some nearest-neighbor outputs are noisy due corpus artifacts.
  \item Next steps:
  \begin{itemize}
    \item curated Azerbaijani analogy benchmark
    \item better class balancing / reweighting
    \item transformer baselines for Task 5
  \end{itemize}
\end{itemize}
\end{frame}

\begin{frame}{Conclusion}
\small
\begin{itemize}
  \item All Project 3 tasks were implemented with reproducible outputs.
  \item Word2Vec and GloVe both learned useful structure with different strengths.
  \item Best downstream model: \texttt{lstm\_\_pmi\_svd} (\(\fmacro=0.5684\)).
  \item UI enables interactive model/result exploration for presentation and grading.
\end{itemize}

\vspace{0.7em}
\centering
\Large Thank you! \\[0.3em]
\normalsize Questions?
\end{frame}

\end{document}
